\begin{question}
The \emph{region labelling problem} is a problem that occurs in image
analysis.  

Assume that each pixel in an image is represented by a boolean: $true$
indicates that the pixel is lit, and $false$ represents that the pixel is
unlit.  The states of the pixels are stored in an array
%
\begin{scala}
  val a = new Array[Array[Boolean]](N,N);
\end{scala}

Let the \emph{neighbours} of a pixel be the (up to four) pixels immediately
above, below, left or right of it.  (Note that the top and bottom rows are not
considered adjacent, and the left and right columns are not considered
adjacent.)  Two pixels are in the same \emph{region} if (and only if) there is
a path from one to the other, where all pixels in the path are lit, and
successive pixels in the path are neighbours.  For example in the image below,
there are three regions.
%
\begin{center}
\includegraphics[width=2cm]{regions.eps}
\end{center}

The aim of this question is to identify the distinct regions and to allocate
each a unique label.  Each label will be of type:
%
\begin{scala}
  type Coord = Tuple2[Int,Int];
\end{scala}
%
More precisely, the label for each region will be the coordinates of the
left-most pixel in the top row of the region, i.e., the one that is smallest
according to the lexicographic ordering.
%
%% \begin{scala}
%%   def smaller(p1:Coord, p2:Coord) : Boolean = {
%%     val i1=p1._1; val j1 = p1._2; 
%%     val i2=p2._1; val j2 = p2._2;
%%     (i1<i2 || i1==i1 && j1<j2)
%%   }
%% \end{scala}
%
An algorithm to calculate this is initially to label each pixel with its own
coordinates; and then to repeatedly relabel each pixel with the smallest label
of itself or its lit neighbours.
%
%% \begin{scala}
%%   for((i1, j1) <- neighbours(i,j))
%%     if(a(i1)(j1) && smaller(label(i1)(j1), label(i)(j)))
%%       label(i)(j) = label(i1)(j1)
%% \end{scala}
%% %
%% where \SCALA{neighbours(i,j) : List[Coord]} gives the neighbours of
%% \SCALA{(i,j)} in some order.

%%%%%

\begin{enumerate}
\item
Write a concurrent program to solve the region labelling problem.  The program
should use \SCALA{p} worker processes (you may assume \SCALA{N\%p = 0}); each
process should work on \SCALA{height = N/p} rows.  (Each process may use an
\SCALA{N} by \SCALA{N} array to store the labels, but will only work on a part
of the array.)  

The worker processes should communicate by message passing only: they may use
no shared variables except they may all read the array~\SCALA{a}.  They may
use a barrier synchronisation to determine if they can terminate.
\marks{20}

%% You should give the definition of a worker, and state carefully what you
%% assume about any channels you use, but do not need to give an explicit
%% definition of how the system is constructed. 

%%%%%

\item
Now suppose a barrier synchronisation operation is not available.  Describe
the changes necessary to your code.
\marks{10}
\end{enumerate}

Your answers to both parts should include suitable test code.

{\bf(Total marks: 30)}
\end{question}

%%%%%%%%%%%%%%%%%%%%%%%%%%%%%%%%%%%%%%%%%%%%%%%%%%%%%%%

\begin{answer}
\Footnotesize
{\bf Part a.}
%
\begin{scala}
// Region labelling problem

import ox.CSO._

object RegionLabelling{ 
  val N = 12; // size of array
  val a = new Array[Array[Boolean]](N,N);

  // Initialise array
  def Init = 
    for(i <- 0 until N; j <- 0 until N) 
      a(i)(j) = (i%3!=1 && j%3!=1)
//    a(i)(j) = 
//      (i!=0 && j%2!=0 || i==1 && j%4==0 || i==N-1 && j%4==2)

  val p = 4; // number of workers
  assert (N%p==0);
  val height = N/p; // height of one strip

  type Coord = Tuple2[Int,Int]; // Cartesian coordinates

  def smaller(p1:Coord, p2:Coord) : Boolean = {
    val i1=p1._1; val j1 = p1._2; val i2=p2._1; val j2 = p2._2;
    (i1<i2 || i1==i1 && j1<j2)
  }

  // find neighbours of p
  def neighbours(p:Coord) : List[Coord] = { 
    val i=p._1; val j=p._2; var ns = Nil : List[Coord];
    if(i>0) ns = (i-1,j)::ns; if(i<N-1) ns = (i+1,j)::ns;
    if(j>0) ns = (i,j-1)::ns; if(j<N-1) ns = (i,j+1)::ns;
    ns;
  }

  // Type of messages, one row of array
  type Msg = Array[Coord]

  // Combining barrier
  def and(b1:Boolean,b2:Boolean) = b1 && b2
  val barrier = new CombiningBarrier(p, true, and);

  def Worker(me: Int, sendUp: ![Msg], sendDown: ![Msg], 
	     receiveUp: ?[Msg], receiveDown: ?[Msg]) 
  = proc("Worker"+me){
    // This worker deals with rows [start .. end)
    val start = me*height; val end = start+height;

    val myLabels = new Array[Array[Coord]](N, N);
    for(i <- 0 until N; j <- 0 until N) 
      if(a(i)(j)) myLabels(i)(j) = (i, j);

    var done = false;

    while(!done){
      var myDone = true;
      // Update coordinates
      for(i <- start until end; j <- 0 until N){
	if(a(i)(j)){
	  for( (i1, j1) <- neighbours(i,j) )
	    if(a(i1)(j1) && 
               smaller(myLabels(i1)(j1), myLabels(i)(j))){
	      myLabels(i)(j) = myLabels(i1)(j1); myDone = false;
	    }
	}
      }
      // Distribute top and bottom rows	  
      if(me!=0) sendUp! myLabels(start).toArray; 
      if(me!=p-1) sendDown! myLabels(end-1).toArray
      // Receive new top and bottom rows
      if(me!=0) myLabels(start-1) = receiveDown?;
      if(me!=p-1) myLabels(end) = receiveUp?;
      done = barrier.sync(myDone);
    }

    // Print results in order of process number, 
    // using down channels as signals
    if(me!=0) receiveDown?; // Wait for previous process
    // Print my rows
    for(i <- start until end){
      print(i+"\t");
      for(j <- 0 until N) print(myLabels(i)(j)+"\t");
      println;
    }
    if(me!=p-1) sendDown!myLabels(height); // Signal to next row
  }

  // Put it together
  val ups = Buf[Msg](1,p); val downs = Buf[Msg](1,p); 
    // indexed by senders

  def System = 
    || ( for(i <- 0 until p) yield 
        Worker(i, ups(i), downs(i), 
               ups((i+1)%p), downs((i+p-1)%p))
      )

  def main(args:Array[String]){
    Init; System();
  }
}
\end{scala}






%% (a) Well assume buffered channels with capacity~1, connected as the names
%%   suggest.
%% %
%% \begin{scala}
%% // Type of messages, one row of array
%% type Msg = Array[Coord]

%% // Combining barrier
%% def and(b1:Boolean,b2:Boolean) = b1 && b2
%% val barrier = new CombiningBarrier(p, true, and);

%% def Worker(me: Int, sendUp: ![Msg], sendDown: ![Msg], 
%%            receiveUp: ?[Msg], receiveDown: ?[Msg]) 
%% = proc("Worker"+me){
%%   // This worker deals with rows [start .. end)
%%   val start = me*height; val end = start+height;

%%   val myLabels = new Array[Array[Coord]](N, N);
%%   for(i <- 0 until N; j <- 0 until N) 
%%     if(a(i)(j)) myLabels(i)(j) = (i, j);

%%   var done = false;

%%   while(!done){
%%     var myDone = true;
%%     // Update coordinates
%%     for(i <- start until end; j <- 0 until N){
%%       if(a(i)(j)){
%%         for( (i1, j1) <- neighbours(i,j) )
%%           if(a(i1)(j1) && 
%%              smaller(myLabels(i1)(j1), myLabels(i)(j))){
%%             myLabels(i)(j) = myLabels(i1)(j1); myDone = false;
%%           }
%%       }
%%     }
%%     // Distribute top and bottom rows   
%%     if(me!=0) sendUp! myLabels(start).toArray; 
%%     if(me!=p-1) sendDown! myLabels(end-1).toArray;
%%     // Receive new top and bottom rows
%%     if(me!=0) myLabels(start-1) = receiveDown?;
%%     if(me!=p-1) myLabels(end) = receiveUp?;
%%     done = barrier.sync(myDone);
%%   }
%% }
%% \end{scala}
%% %
%% [The call to \SCALA{toArray} creates a fresh copy of the array, to avoid
%%   sharing.] 

%%%%%

(b) We piggy-back the booleans indicating whether processes can terminate onto
the data.  Define the type of messages by
%
\begin{scala}
type Msg = Tuple2[Array[Coord], Boolean];
\end{scala}
%
And replace the data-passing and global synchronisation by
%
\begin{scala}
if(me==0){
  // Send my bottom row
  sendDown! (myLabels(end-1).toArray, myDone); 
  val (newBelow, b2) = receiveUp?;  // Receive new row below
  myLabels(end) = newBelow; done = b2;
}
else if(me==p-1){
  val (newAbove, b1) = receiveDown?; // Receive new row above
  myLabels(start-1) = newAbove; done = b1 && myDone;
  sendUp! (myLabels(start).toArray, done); // Send my top row
}
else{
  val (newAbove, b1) = receiveDown?; // Receive new row above
  myLabels(start-1) = newAbove;
  // Send my bottom row
  sendDown! (myLabels(end-1).toArray, myDone && b1); 
  val (newBelow, b2) = receiveUp?;  // Receive new row below
  myLabels(end) = newBelow; done = b2;
  sendUp! (myLabels(start).toArray, b2); // Send my top row
}
\end{scala}
%
The bottom rows are passed down, during which the conjunction of the
\SCALA{myDone} variables is calculated; then the top rows are passed up with
the final value of \SCALA{done}.
\end{answer}
